\section*{अध्याय \ref{ch:b1:e-ph-diagrams} --- विभव--pH आरेख}

\begin{solution}{exo:b1:e-ph-diagrams:1}
मैंगनीज़: \ce{Mn} (0) $<$ \ce{Mn^2+}, \ce{Mn(OH)2} ($+$II) $<$ \ce{MnO2}
($+$IV) $<$ \ce{MnO4-} ($+$VII)। क्लोरीन: \ce{Cl-} ($-$I) $<$ \ce{Cl2} (0)
$<$ \ce{HClO}, \ce{ClO-} ($+$I) $<$ \ce{ClO3-} ($+$V)।
\end{solution}

\begin{solution}{exo:b1:e-ph-diagrams:2}
$-0.059 \times 3/1 = -0.177$; $+0.059 \times 2/2 = +0.059$;
$-0.059 \times 8/5 = -0.094$; $-0.059 \times 2/2 = -0.059$ V प्रति pH इकाई।
कॉपर युग्म के लिए प्रोटॉन अपचित रूप की ओर मुक्त होते हैं
($m = -2$): विभव pH के साथ बढ़ता है।
\end{solution}

\begin{solution}{exo:b1:e-ph-diagrams:3}
\ce{H+}/\ce{H2}, pH 0 पर 0~V से गुज़रती है और $\mathrm{pH} \geqslant 0$ के लिए
0.50~V तक कभी नहीं पहुँचती। \ce{O2}/\ce{H2O},
$\mathrm{pH} = 0.73/0.059 = 12.4$ पर 0.50~V तक पहुँचती है और 0~V तक केवल pH 20.8 पर, जो
पैमाने से बाहर है। पट्टी हर pH पर \qty{1.23}{V} चौड़ी है, pH 7 पर भी
($-0.41$ से $+0.82$~V तक)।
\end{solution}

\begin{solution}{exo:b1:e-ph-diagrams:4}
(1, 1.0~V): \ce{Fe^3+}। (4, 0): रेखा \ce{Fe(OH)3}/\ce{Fe^2+}
$1.09 - 0.71 = 0.38$~V पर है और बिंदु उसके नीचे, $-0.47$~V के ऊपर है:
\ce{Fe^2+}। (10, $-0.4$~V): $-0.07 - 0.59 = -0.66$ और $0.28 -
0.59 = -0.31$~V के बीच: \ce{Fe(OH)2}। (10, 0.5~V): \ce{Fe(OH)3}।
\end{solution}

\begin{solution}{exo:b1:e-ph-diagrams:5}
\ce{Fe^2+}/\ce{Fe}: $-0.41 + 0.030 \times (-6) = -0.59$~V।
\ce{Fe^3+}/\ce{Fe(OH)3}: $14.00 - \frac13(38.55 - 6) = 3.15$। घुले आयनों के
क्षेत्र (संक्षारण) बढ़ते हैं: \ce{Fe^2+} नीचे $-0.59$~V तक फैल जाता है
और \ce{Fe^3+}, pH 3.15 तक; धातु और हाइड्रॉक्साइड क्षेत्र खोते हैं।
\end{solution}

\begin{solution}{exo:b1:e-ph-diagrams:6}
\ce{Fe(OH)3 + 3H+ + e- -> Fe^2+ + 3H2O}: $E = E^\circ + 0.059\log\frac{[\ce{H+}]^3}{[\ce{Fe^2+}]}
= (E^\circ - 0.059\log c) - 0.177\,\mathrm{pH}$। pH 1.82 पर रेखा
\ce{Fe^3+}/\ce{Fe^2+} से 0.77~V पर मिलती है, इसलिए स्थिरांक $0.77 + 0.177 \times 1.82
= 1.09$~V है।
\end{solution}

\begin{solution}{exo:b1:e-ph-diagrams:7}
$E^\circ(\ce{Cu+}/\ce{Cu}) = 0.52 > E^\circ(\ce{Cu^2+}/\ce{Cu+}) = 0.16$:
\cref{prop:b1:e-ph-diagrams:disproportionation-criterion} के अनुसार \ce{Cu+}
असमानुपातित होता है। एकल सीमा: $E = 0.34 + 0.030\log 0.010 =
0.28$~V।
\end{solution}

\begin{solution}{exo:b1:e-ph-diagrams:8}
pH 0 पर कॉपर का क्षेत्र (0.28~V से नीचे) पानी के क्षेत्र से अतिव्याप्त होता है और
हाइड्रोजन की रेखा के ऊपर है: वायु के बिना हाइड्रोक्लोरिक अम्ल से कोई अभिक्रिया
नहीं। नाइट्रिक अम्ल: 0.96~V पर \ce{NO3-}/\ce{NO} कॉपर के
क्षेत्र से ऊपर है, जो ऑक्सीकृत हो जाता है:
\ce{3Cu + 2NO3- + 8H+ -> 3Cu^2+ + 2NO + 4H2O}।
% ledger: dfg:NO3-_ao, dfg:NO_g (E0 computed in figdata/bachelor-1/standard-potentials.py)
\end{solution}

\begin{solution}{exo:b1:e-ph-diagrams:9}
pH 3 पर रेखा \ce{O2}/\ce{H2O}, $1.23 - 0.18 = 1.05$~V पर है, \ce{Fe^2+} के
पूरे क्षेत्र से ऊपर: ऑक्सीजन आयरन(II) को ऑक्सीकृत करती है। pH 1.82 से ऊपर
बना आयरन(III) पीले-भूरे हाइड्रॉक्साइड के रूप में अवक्षेपित होता है:
\ce{4Fe^2+ + O2 + 10H2O -> 4Fe(OH)3 + 8H+}।
\end{solution}

\begin{solution}{exo:b1:e-ph-diagrams:10}
ऊर्ध्वाधर सीमाएँ: $14.00 - \frac12(16.38 - 2) = 6.81$; और,
\ce{Zn(OH)2 + 2OH- <=> Zn(OH)4^2-} ($K = 10^{-1.93}$) से,
$[\ce{Zn(OH)4^2-}] = 10^{-2}$ के साथ, $[\ce{OH-}]^2 = 10^{-0.07}$, pH 13.97।
\ce{Zn^2+}/\ce{Zn}: $-0.76 + 0.030 \times (-2) = -0.82$~V। \ce{Zn(OH)2}/\ce{Zn}:
प्रवणता $-0.059$, $(6.81, -0.82)$ से होकर: $E = -0.42 - 0.059\,\mathrm{pH}$।
\ce{Zn(OH)4^2- + 4H+ + 2e- -> Zn + 4H2O}: प्रवणता $-0.118$,
$(13.97, -1.24)$ से होकर: $E = 0.41 - 0.118\,\mathrm{pH}$।
\end{solution}

\begin{solution}{exo:b1:e-ph-diagrams:11}
आयरन का क्षेत्र हर pH पर हाइड्रोजन रेखा के नीचे है। pH 2: आयरन
हाइड्रोजन की मुक्ति के साथ \ce{Fe^2+} में ऑक्सीकृत होता है (संक्षारण)। pH 7:
अभिक्रिया अब भी संभव है, पर हाइड्रोजन रेखा ($-0.41$~V)
\ce{Fe^2+}/\ce{Fe} के पास है और अभिक्रिया बहुत धीमी है। pH 13: आयरन
\ce{Fe(OH)2} में ऑक्सीकृत होता है, जो \omterm{def:b1:e-ph-diagrams:corrosion-domains}{निष्क्रियण क्षेत्र} है: एक परत बनती है और
कील चमकती रहती है। घुली ऑक्सीजन के साथ रेखा \ce{O2}/\ce{H2O},
जो बहुत ऊपर है, हर pH पर आयरन को आयरन(III) में ऑक्सीकृत करती है: जंग, \ce{Fe(OH)3}
और उसके निर्जलित ऑक्साइड।
\end{solution}

\begin{solution}{exo:b1:e-ph-diagrams:12}
ज़िंक का क्षेत्र नीचे है: इस्पात के संपर्क में वह पहले
ऑक्सीकृत होता है, pH 8 पर \ce{2Zn + O2 + 2H2O -> 2Zn(OH)2}, और उससे
मुक्त इलेक्ट्रॉन इस्पात की सतह पर ऑक्सीजन को अपचित करते हैं, जो धात्विक बनी रहती है।
मैग्नीशियम ($E^\circ = -2.36$~V) और भी नीचे है और वह भी काम करता है; कॉपर
(0.34~V) आयरन से ऊपर है: कॉपर से जुड़ा होने पर आयरन ही
ऑक्सीकृत होने वाली धातु होता, और अकेले की तुलना में अधिक तेज़ी से।
\end{solution}

\begin{solution}{pb:b1:e-ph-diagrams:1}
\textbf{1.} $-$I, 0, $+$I, $+$I।
\textbf{2.} सबसे नीचे \ce{Cl-}, फिर \ce{Cl2}, फिर सबसे ऊपर \ce{HClO} और \ce{ClO-}।
\textbf{3.} \ce{HClO} और \ce{ClO-}, एक \omterm{def:b1:predominant-reaction:bronsted}{अम्ल--क्षारक युग्म}।
\textbf{4.} \ce{Cl2 + 2e- -> 2Cl-}।
\textbf{5.} \ce{2HClO + 2H+ + 2e- -> Cl2 + 2H2O}।
\textbf{6.} \ce{ClO- + H2O + 2e- -> Cl- + 2OH-}।
\textbf{7.} pH 7.55 से नीचे \ce{HClO}, ऊपर \ce{ClO-}।
\textbf{8.} कोई इलेक्ट्रॉन विनिमित नहीं होता; सीमा केवल $K_a$ से तय होती है।
\textbf{9.} $1/(1 + 10^{7.4 - 7.55}) = 0.59$: \qty{59}{\percent}।
\textbf{10.} $1/(1 + 10^{0.45}) = 0.26$: सक्रिय क्लोरीन का तीन-चौथाई भाग
कम प्रभावी रोगाणुनाशक \ce{ClO-} होता।
\textbf{11.} \ce{ClO-}।
\textbf{12.} $E = 1.396 + 0.0296\log\frac{[\ce{Cl2}]}{[\ce{Cl-}]^2} = 1.396
+ 0.0296\log\frac{0.005}{10^{-4}} = 1.396 + 0.050 = 1.446$~V; \omterm{def:b1:nernst:couple}{अर्ध-समीकरण} में
कोई \ce{H+} नहीं।
\textbf{13.} $E = 1.594 + 0.0296\log\frac{[\ce{HClO}]^2[\ce{H+}]^2}{[\ce{Cl2}]}
= 1.594 + 0.0296\log\frac{10^{-4}}{0.005} - 0.0592\,\mathrm{pH} = 1.594 -
0.050 - 0.0592\,\mathrm{pH}$।
\textbf{14.} बराबर रखने पर $0.0592\,\mathrm{pH} = 1.594 - 1.396 - 2 \times 0.0503
= 0.0974$, $\mathrm{pH} = 1.646$।
\textbf{15.} \ce{HClO + H+ + 2e- -> Cl- + H2O}: दो इलेक्ट्रॉनों पर एक प्रोटॉन,
प्रवणता $-0.0592/2 = -0.0296$।
\textbf{16.} दोनों \omterm{def:b1:nernst:couple}{अर्ध-समीकरणों} को जोड़ने पर (हर एक में प्रति क्लोरीन परमाणु एक
इलेक्ट्रॉन): $2E^\circ = 1.594 + 1.396$, $E^\circ(\ce{HClO}/\ce{Cl-}) =
1.495$~V; सीमा $E = 1.495 - 0.0296\,\mathrm{pH}$।
\textbf{17.} \ce{ClO- + 2H+ + 2e- -> Cl- + H2O}: $E = 1.718 - 0.0592\,\mathrm{pH}$
(स्थिरांक सातत्य के लिए चुना गया: $1.495 - 0.0296 \times 7.55 = 1.272 = 1.718
- 0.0592 \times 7.55$)।
\textbf{18.} रेखाओं के नीचे \ce{Cl-}; 1.446~V के ऊपर,
$\mathrm{pH} < 1.65$ पर एक छोटे त्रिभुज में \ce{Cl2}; pH 7.55 तक प्रवणता
$-0.0296$ वाली रेखा के ऊपर \ce{HClO}, उसके बाद \ce{ClO-}। रेखा \ce{O2}/\ce{H2O},
1.23 से 0.40~V तक, इन सभी सीमाओं के नीचे है (अध्याय की स्क्रिप्ट से परिकलित
आरेख का आकार भी यही है)।
\textbf{19.} नहीं: इसका क्षेत्र ऑक्सीजन की रेखा के ऊपर है, इसलिए यह
पानी को ऑक्सीकृत कर सकता है, पर अभिक्रिया धीमी है। तरणताल अपनी क्लोरीन मुख्यतः
इसलिए खोता है कि वह तैराकों और हवा द्वारा लाई गई चीज़ों को ऑक्सीकृत करती है, और
धूप के कारण; उसे फिर से भरना पड़ता है।
\textbf{20.} वहाँ \ce{Cl2} का कोई क्षेत्र नहीं है: यह असमानुपातित होता है,
\ce{Cl2 + H2O -> HClO + Cl- + H+}; क्षारक में
\ce{Cl2 + 2OH- -> ClO- + Cl- + H2O}।
\textbf{21.} सोडियम हाइड्रॉक्साइड विलयन में क्लोरीन प्रवाहित करके (प्रश्न 20 की
अभिक्रिया)।
\textbf{22.} \ce{HClO + Cl- + H+ -> Cl2 + H2O}, एक \omterm{def:b1:e-ph-diagrams:disproportionation}{समानुपातन}।
\textbf{23.} pH 1.6 से नीचे हाइपोक्लोराइट और क्लोराइड समानुपातित होते हैं: क्लोरीन बनती है,
और अपनी \omterm{def:b1:precipitation:solubility}{विलेयता} से अधिक होने पर विषैली गैस के रूप में निकल जाती है।
\textbf{24.} $c = \qty{0.010}{mol/L}$ पर pH 4,
\ce{Cl2} के क्षेत्र से बाहर है। पर सीमा $c$ पर निर्भर करती है: किसी व्यापक $c$ के लिए
प्रश्न 12--14 दोहराने पर $\mathrm{pH} = 3.34 + \log(2c)$ मिलता है, जो
$c = \qty{1}{mol/L}$ के लिए, जो ब्लीच की सांद्रता की कोटि है, 3.6 तक पहुँचता है।
pH 4 के पास सांद्र उत्पादों को मिलाना भी सुरक्षित नहीं है।
\textbf{25.} \textbf{pH 1.6}: इसके ऊपर, $c = \qty{0.010}{mol/L}$ पर,
घुली क्लोरीन हाइपोक्लोरस अम्ल और क्लोराइड में असमानुपातित होती है।
\end{solution}
